\documentclass[conference]{IEEEtran}

\usepackage[utf8]{inputenc}
\usepackage[T1]{fontenc}
\usepackage{amsmath,amssymb}
\usepackage{graphicx}
\usepackage{booktabs}
\usepackage{siunitx}

\title{Modelado, Simulaci\'on y Control de una Plataforma Antivibratoria Activa}

\author{
\IEEEauthorblockN{Nombre 1, Nombre 2}
\IEEEauthorblockA{
Universidad de los Andes, Bogot\'a, Colombia\\
correo1@uniandes.edu.co, correo2@uniandes.edu.co}
}

\begin{document}

\maketitle

\begin{abstract}
En este trabajo se estudia una plataforma antivibratoria activa destinada al aislamiento de un equipo sensible montado sobre una base con vibraci\'on. Se construye un modelo realista con no linealidades f\'isicas y un modelo matem\'atico aproximado de tercer orden para dise\~nar un controlador mediante control por modelo interno (IMC). Se comparan ambos modelos en lazo abierto y en lazo cerrado, evaluando seguimiento de referencia, rechazo a perturbaciones y esfuerzo de control.
\end{abstract}

\section{Introducci\'on}
La reducci\'on de vibraciones transmitidas a instrumentos sensibles es un problema cl\'asico de control. En esta actividad se plantea una plataforma antivibratoria activa en la que una base inferior est\'a sometida a vibraciones externas y una plataforma superior debe permanecer cercana a una referencia deseada. El inter\'es del proyecto radica en la diferencia entre un modelo realista, que incluye no linealidades y restricciones f\'isicas, y un modelo lineal aproximado usado para dise\~nar el controlador.

\section{Modelo realista}
La plataforma se modela como una masa $m$ conectada a la base por un resorte y un amortiguador. Adem\'as, un actuador aplica una fuerza de control $F_a$ y el modelo realista incluye rigidez c\'ubica, fricci\'on tipo Coulomb suavizada, saturaci\'on del actuador y topes mec\'anicos.

Definiendo $x(t)$ como la posici\'on absoluta de la plataforma, $z(t)$ como el movimiento de la base y $u(t)$ como la se\~nal de control, el modelo realista se describe por
\begin{equation}
m\ddot{x}=F_a-F_{susp}-F_{stop},
\end{equation}
con
\begin{equation}
F_{susp}=k(x-z)+c(\dot{x}-\dot{z})+k_3(x-z)^3+F_c\tanh\left(\frac{\dot{x}-\dot{z}}{v_\varepsilon}\right),
\end{equation}
y din\'amica del actuador
\begin{equation}
\tau_a \dot{F}_a + F_a = K_a \,\mathrm{sat}(u).
\end{equation}

\section{Modelo aproximado y estabilidad}
Al despreciar rigidez c\'ubica, fricci\'on no lineal, topes mec\'anicos y saturaci\'on, y linealizar alrededor del equilibrio $x^\star=\dot{x}^\star=z^\star=F_a^\star=u^\star=0$, se obtiene
\begin{equation}
m\ddot{x}+c\dot{x}+kx=F_a,
\end{equation}
\begin{equation}
\tau_a \dot{F}_a+F_a=K_a u.
\end{equation}
La funci\'on de transferencia entre entrada de control y posici\'on de la plataforma es
\begin{equation}
G(s)=\frac{X(s)}{U(s)}=\frac{K_a}{(\tau_a s+1)(m s^2+c s+k)}.
\end{equation}

Con los valores $m=\SI{12}{kg}$, $k=\SI{1800}{N/m}$, $c=\SI{95}{N\,s/m}$, $\tau_a=\SI{0.04}{s}$ y $K_a=\SI{220}{N/V}$, el denominador es
\begin{equation}
0.48s^3+15.8s^2+167s+1800.
\end{equation}
Los polos asociados tienen parte real negativa, por lo que la planta aproximada es estable en lazo abierto.

\section{Dise\~no del controlador}
Se emplea control por modelo interno con filtro
\begin{equation}
F(s)=\frac{1}{(\lambda s+1)^3},
\end{equation}
seleccionando $\lambda=\SI{0.10}{s}$. Como la planta es estable e invertible, el controlador interno es
\begin{equation}
Q(s)=G^{-1}(s)F(s).
\end{equation}
El controlador equivalente en realimentaci\'on queda
\begin{equation}
G_c(s)=\frac{(\tau_a s+1)(m s^2+c s+k)}
{K_a \lambda s(\lambda^2 s^2+3\lambda s+3)}.
\end{equation}
Para la planta nominal, la funci\'on de transferencia en lazo cerrado es
\begin{equation}
T(s)=\frac{1}{(\lambda s+1)^3},
\end{equation}
lo cual anticipa una respuesta monotona y sin sobreimpulso nominal.

\section{Escenarios de simulaci\'on}
Se comparan el modelo lineal y el modelo realista bajo dos pruebas:
\begin{enumerate}
\item Respuesta en lazo abierto ante un escal\'on de \SI{0.05}{V}.
\item Respuesta en lazo cerrado con referencia variable:
\[
0 \rightarrow \SI{4}{mm} \rightarrow \SI{-2.5}{mm} \rightarrow \SI{3}{mm},
\]
con una perturbaci\'on de base sinusoidal aplicada a partir de \SI{6}{s}.
\end{enumerate}

\section{Resultados esperados}
Al ejecutar los scripts en MATLAB 2025a se obtienen resultados cercanos a los siguientes:
\begin{table}[!t]
\caption{Resumen de resultados nominales}
\label{tab:resumen}
\centering
\scriptsize
\begin{tabular}{lcc}
\toprule
Caso & $M_p$ & $t_{ss}$ \\
\midrule
Lazo abierto lineal & 29.96\% & --- \\
Lazo abierto realista & 18.07\% & --- \\
Lazo cerrado lineal & 0\% & 0.90 s \\
Lazo cerrado realista & 5.9\% & 2.95 s \\
\bottomrule
\end{tabular}
\end{table}

Durante la ventana de perturbaci\'on, el error RMS esperado es del orden de mil\'imetros, lo cual permite discutir de forma clara el rechazo a vibraciones de base y la diferencia entre el modelo lineal y el realista.

\section{Implementaci\'on}
Se entrega una carpeta con:
\begin{itemize}
\item scripts MATLAB para an\'alisis y simulaci\'on;
\item modelos Simulink generables autom\'aticamente;
\item un documento detallado del proyecto;
\item el presente borrador en formato IEEE.
\end{itemize}

\section{Conclusiones}
La plataforma antivibratoria activa propuesta cumple las condiciones de la actividad: es un sistema SISO, posee una din\'amica aproximada de tercer orden y permite comparar un modelo realista con uno lineal para dise\~nar un controlador. El m\'etodo IMC produce un controlador coherente con el temario del curso y permite estudiar tanto seguimiento de referencia como rechazo a perturbaciones.

\end{document}
